\documentclass[11pt]{article}
% ===========================================================================
%  Shared preamble for the AP Calculus AB exam-preparation worksheets.
%  Input right after \documentclass. Files beginning with "_" are not compiled
%  on their own. Navy and gold, to match the course's lesson pages.
% ===========================================================================
\usepackage[letterpaper,top=2.3cm,bottom=2.2cm,left=1.7cm,right=1.7cm,headheight=28pt,footskip=30pt]{geometry}
\usepackage{amsmath,amssymb}
\usepackage[T1]{fontenc}
\usepackage[utf8]{inputenc}
\usepackage{lmodern}
\usepackage{xcolor}
\usepackage[most]{tcolorbox}
\usepackage{enumitem}
\usepackage{fancyhdr}
\usepackage{array}
\usepackage{tikz}
\usepackage{pgfplots}
\pgfplotsset{compat=1.18}

\definecolor{navy}{HTML}{161E54}
\definecolor{ink}{HTML}{27357E}
\definecolor{gold}{HTML}{B98A00}
\definecolor{goldbg}{HTML}{FFFDF3}
\definecolor{inkbg}{HTML}{F6F7FD}
\definecolor{rule}{HTML}{C9CFF0}
\definecolor{okgreen}{HTML}{1B7A44}
\definecolor{okbg}{HTML}{F1F8F3}

\pagestyle{fancy}
\fancyhf{}
\fancyhead[L]{\footnotesize Name: \underline{\hspace{4cm}}}
\fancyhead[C]{\textbf{\color{navy}Integration Academy}\\[-2pt]{\scriptsize Calculus AB \textperiodcentered{} AP\textsuperscript{\textregistered} Exam Preparation}}
\fancyfoot[L]{\tiny\color{black!55}\parbox[t]{0.78\textwidth}{AP\textsuperscript{\textregistered} is a registered trademark of the College Board, which was not involved in the production of, and does not endorse, this material. \textcopyright{} Integration Academy.}}
\fancyfoot[R]{\scriptsize\thepage}
\renewcommand{\headrulewidth}{1.2pt}
\renewcommand{\headrule}{\hbox to\headwidth{\color{gold}\leaders\hrule height \headrulewidth\hfill}}
\renewcommand{\footrulewidth}{0pt}
\setlength{\parindent}{0pt}
\setlength{\parskip}{4pt}

% Title block: \wstitle{code}{title}{unit line}
\newcommand{\wstitle}[3]{\begin{center}{\small\bfseries\color{gold}WORKSHEET #1}\\[3pt]{\LARGE\bfseries\color{navy}#2}\\[3pt]{\small\color{black!60}#3}\end{center}\vspace{2pt}}
% Section heading
\newcommand{\sech}[1]{\par\vspace{10pt}{\large\bfseries\color{navy}#1}\par\nobreak\vspace{2pt}{\color{rule}\hrule height0.8pt}\vspace{7pt}}
% A level tag printed beside a group of questions
\newcommand{\level}[1]{\par\vspace{6pt}{\footnotesize\bfseries\color{gold}\MakeUppercase{#1}}\par\nobreak\vspace{2pt}}

\newtcolorbox{ideabox}{colframe=ink,colback=inkbg,coltitle=white,colbacktitle=ink,fonttitle=\bfseries\small,title={Key ideas},arc=2pt,boxrule=0.8pt,left=8pt,right=8pt}
\newtcolorbox{example}[1]{colframe=rule,colback=white,coltitle=navy,colbacktitle=inkbg,fonttitle=\bfseries\small,title={#1},arc=2pt,boxrule=0.8pt,borderline west={3pt}{0pt}{ink},left=10pt,right=8pt}
\newtcolorbox{frbox}[1]{colframe=gold,colback=goldbg,coltitle=white,colbacktitle=gold,fonttitle=\bfseries\small,title={#1},arc=2pt,boxrule=0.8pt,breakable,left=8pt,right=8pt}
\newtcolorbox{solbox}{colframe=okgreen,colback=okbg,coltitle=white,colbacktitle=okgreen,fonttitle=\bfseries\small,title={Worked solutions},arc=2pt,boxrule=0.8pt,breakable,left=8pt,right=8pt}
\newtcolorbox{scorebox}[1]{colframe=navy,colback=white,coltitle=white,colbacktitle=navy,fonttitle=\bfseries\small,title={#1},arc=2pt,boxrule=0.8pt,breakable,left=8pt,right=8pt}

\newcommand{\soln}{\par\smallskip\textbf{\color{ink}Solution.}\ }
\newcommand{\result}[1]{\par\smallskip\textbf{#1}}
% blank working space of a given height
\newcommand{\work}[1]{\par\vspace{#1}}
% a point value in the margin of a scoring guide
\newcommand{\pts}[1]{\hfill{\footnotesize\bfseries\color{gold}[#1]}}

% Coordinate grid for graph questions: \begin{apgraph}{xmin}{xmax}{ymin}{ymax} ... \end{apgraph}
\newenvironment{apgraph}[4]{\begin{center}\begin{tikzpicture}\begin{axis}[
  axis lines=middle, xmin=#1, xmax=#2, ymin=#3, ymax=#4,
  grid=both, grid style={black!12}, xtick distance=1, ytick distance=1,
  ticklabel style={font=\scriptsize}, xlabel={$x$}, ylabel={$y$},
  width=10cm, height=7.2cm, samples=80, clip=false,
  every axis plot/.append style={line width=1.1pt, color=ink}]}{\end{axis}\end{tikzpicture}\end{center}}
\newcommand{\closeddot}[2]{\addplot[only marks, mark=*, mark size=2.2pt, color=ink] coordinates {(#1,#2)};}
\newcommand{\opendot}[2]{\addplot[only marks, mark=*, mark size=2.2pt, color=ink, mark options={fill=white}] coordinates {(#1,#2)};}

% Multiple-choice options: four across, or stacked when the options are long.
\newcommand{\opts}[4]{\par\vspace{3pt}\begin{tabular}{@{}p{0.235\linewidth}p{0.235\linewidth}p{0.235\linewidth}p{0.235\linewidth}@{}}(A)\ #1 & (B)\ #2 & (C)\ #3 & (D)\ #4\end{tabular}\par\vspace{9pt}}
\newcommand{\optsv}[4]{\par\vspace{3pt}\begin{tabular}{@{}l@{\ \ }p{0.9\linewidth}@{}}(A) & #1\\(B) & #2\\(C) & #3\\(D) & #4\end{tabular}\par\vspace{9pt}}

\fancyhead[R]{\footnotesize\color{black!60}Unit 1 free response: scoring guide}
\begin{document}

\wstitle{1.12}{Unit 1 Free-Response Practice: Scoring Guide}{Unit 1: Limits and Continuity \textperiodcentered{} 18 points}

\begin{scorebox}{Question 1 \textperiodcentered{} 9 points}
\textbf{(a)} For $x<2$, $\dfrac{x^2-x-2}{x-2}=\dfrac{(x-2)(x+1)}{x-2}=x+1$, so $\displaystyle\lim_{x\to 2^-}f(x)=3$. \pts{1 point: left-hand limit with supporting work}

$\displaystyle\lim_{x\to 2^+}f(x)=\sqrt{2+7}=3$. \pts{1 point: right-hand limit}

\medskip\textbf{(b)} $k=3$. $f$ is continuous at $x=2$ when $\displaystyle\lim_{x\to 2^-}f(x)=\displaystyle\lim_{x\to 2^+}f(x)=f(2)$, and both one-sided limits equal 3. \pts{1 point: $k=3$; 1 point: justification using the limits}

\medskip\textbf{(c)} A removable discontinuity, because $\displaystyle\lim_{x\to 2}f(x)=3$ exists but $f(2)=5\neq3$. \pts{1 point: type with reason}

\medskip\textbf{(d)} For large $x$, $f(x)=\sqrt{x+7}$, so $\displaystyle\lim_{x\to \infty}\dfrac{\sqrt{x+7}}{\sqrt x}=\displaystyle\lim_{x\to \infty}\sqrt{1+\dfrac7x}=1$. \pts{1 point: uses the correct piece; 1 point: answer}

\medskip\textbf{(e)} Yes. With $k=3$, $f$ is continuous on $[2,9]$. $f(2)=3$ and $f(9)=\sqrt{16}=4$, so $f(2)<3.5<f(9)$. By the Intermediate Value Theorem there is a $c$ in $(2,9)$ with $f(c)=3.5$. \pts{1 point: continuity and $f(2)<3.5<f(9)$; 1 point: conclusion naming the theorem}
\end{scorebox}

\begin{scorebox}{Question 2 \textperiodcentered{} 9 points}
\textbf{(a)} $\dfrac{W(6)-W(2)}{6-2}=\dfrac{46-52}{4}=-1.5$ megalitres per day. \pts{1 point: difference quotient and value}

From $t=2$ to $t=6$ days, the amount of water in the reservoir decreased at an average rate of 1.5 megalitres per day. \pts{1 point: interpretation with units and interval}

\medskip\textbf{(b)} $W$ is continuous on $[0,10]$. $50$ lies between $W(0)=40$ and $W(2)=52$, between $W(2)=52$ and $W(6)=46$, and between $W(6)=46$ and $W(10)=61$. By the Intermediate Value Theorem, $W(t)=50$ at least once in each of $(0,2)$, $(2,6)$ and $(6,10)$: at least $3$ times. \pts{1 point: continuity stated; 1 point: 50 between the values on the three intervals; 1 point: answer 3}

\medskip\textbf{(c)} $\displaystyle\lim_{t\to \infty}\left(60-20e^{-0.3t}\right)=60$, since $e^{-0.3t}\to0$. \pts{1 point: limit}

In the long run the model predicts that the amount of water in the reservoir approaches 60 megalitres. \pts{1 point: interpretation}

\medskip\textbf{(d)} With $t=0.01$: $\dfrac{M(0.01)-40}{0.01}=5.991$. With $t=-0.01$: $6.009$. With $t=\pm0.001$: $5.999$ and $6.001$. The limit is approximately $6.000$. \pts{1 point: values on both sides of 0; 1 point: estimate}
\end{scorebox}

\begin{scorebox}{Using your score}
Mark each question out of 9. As a rough guide for a single unit: 14 or more out of 18 is strong; 10 to 13 is on track; below 10, rework the lessons named beside the points you missed (limits: 1.5; continuity: 1.8; discontinuity types: 1.7; limits at infinity: 1.9; Intermediate Value Theorem: 1.10; rates and interpretation: 1.1; tables: 1.3).
\end{scorebox}
\end{document}
